% GitHub Repo and Documentation: https://github.com/celiobjunior/resume-template
% Copyright © 2025 Celio B Junior. All rights reserved.
%
% Licensed under the Apache License, Version 2.0 (the "License");
% you may not use this file except in compliance with the License.
% You may obtain a copy of the License at
%
%     http://www.apache.org/licenses/LICENSE-2.0
%
% This template follows best practices from README.md

% Início do documento LaTeX: define tipo e formato do documento
% a4paper = tamanho A4, 10pt = tamanho base da fonte
\documentclass[a4paper,10pt]{article}

% --- PACOTES: BASE / IDIOMA ---
\usepackage[utf8]{inputenc}
\usepackage[T1]{fontenc}
\usepackage[portuguese]{babel}

% --- PACOTES: LAYOUT / TIPOGRAFIA ---
\usepackage{geometry}
\usepackage{parskip}
\usepackage{microtype}
\usepackage{enumitem}
\usepackage{titlesec}
\usepackage{array}
\usepackage{tabularx}

% --- PACOTES: LINKS / BOOKMARKS ---
\usepackage{hyperref}
\usepackage{bookmark}
\usepackage{xurl}

% --- CONFIGURAÇÃO: MARGENS / ESPAÇAMENTO ---
\geometry{top=1.5cm, bottom=1.5cm, left=1.5cm, right=1.5cm} % Margens da página
\setcounter{secnumdepth}{0}                                 % Remove numeração de seção
\setlist[itemize]{
    leftmargin=0.75em, % Recuo da lista (mais perto da margem da página)
    itemsep=0.2em,     % Espaço entre itens
    topsep=0.25em,     % Espaço antes/depois da lista
    parsep=0em,        % Espaço entre parágrafos no mesmo item
    partopsep=0em      % Espaço extra ao iniciar a lista
}

% Remove números de página e cabeçalhos para visual limpo do currículo
\pagestyle{empty}

% --- CONFIGURAÇÃO: METADADOS / LINKS ---
\hypersetup{
    pdftitle={CV Otávio Araujo},   % Título do PDF
    pdfauthor={Otávio Araujo},     % Autor do PDF
    colorlinks=true,          % Usa cor nos links (sem caixa)
    linkcolor=black,          % Cor de links internos
    urlcolor=black,           % Cor de URLs
    citecolor=black,          % Cor de citações
    bookmarksdepth=2          % Barra lateral: seção e subseção
}

% --- CONFIGURAÇÃO: TÍTULOS ---
\titleformat{\section}
{\Large\bfseries}
{}
{0em}
{}
[\titlerule\vspace{0.5ex}]
\titleformat{\subsection}
{\normalsize\bfseries}
{}
{0em}
{}
\titlespacing*{\section}{0pt}{0.7em}{0.25em}
\titlespacing*{\subsection}{0pt}{0.35em}{0.15em}

% Macro para entradas do currículo com bookmark no PDF
\newcounter{cventry}
\newcommand{\cventry}[4]{%
    \refstepcounter{cventry}%
    \phantomsection%
    \pdfbookmark[2]{#1}{cventry-\thecventry}%
    \noindent\begin{tabularx}{\textwidth}{@{}>{\raggedright\arraybackslash}X >{\raggedleft\arraybackslash}X@{}}
    \textbf{#1} & #2 \\
    \textit{#3} & \textit{#4} \\
    \end{tabularx}
}

% --- INÍCIO DO DOCUMENTO ---
\begin{document}

% --- CABEÇALHO ---
\begin{center}
    {\LARGE \textbf{Otávio Araujo}}
    \\ [0.1cm]
    Belém, Pará
    {\textbullet}
    \href{mailto:otaviodev08@gmail.com}{otaviodev08@gmail.com}
    {\textbullet}
    \href{https://linkedin.com/in/otávio-araujo-77474b1ab}{linkedin.com/in/otávio-araujo-77474b1ab}
    {\textbullet}
    \href{https://github.com/OtavioAraujoS}{github.com/OtavioAraujoS}
\end{center}

% --- SEÇÕES ---

\section{Experiência}
    \cventry{Deskcorp}{Belém, Brasil}{Arquiteto de Soluções}{Julho de 2026 -- Atualmente}
        \begin{itemize}
            \item Desenhei arquiteturas de software e modelei processos de negócio utilizando Bizagi e draw.io, elaborando diagramas de solução, classes e objetos para traduzir requisitos complexos em visões técnicas claras para os times de desenvolvimento.
            \item Atuei como facilitador entre áreas técnicas e de negócios, conduzindo reuniões com times multidisciplinares para levantar requisitos, definir arquiteturas e refinar regras de negócio.
            \item Mapeei e gerenciei riscos operacionais e tecnológicos, desenvolvendo estratégias de mitigação e planos de contingência para aumentar a resiliência, a segurança e a estabilidade das entregas.
            \item Administrei identidades e acessos (IAM) em múltiplas plataformas, garantindo o provisionamento seguro de clientes conforme as políticas de conformidade e segurança da informação.
            \item Gerenciei o ciclo de vida da documentação técnica e dos projetos, estruturando e padronizando os artefatos de entrega para assegurar governança, rastreabilidade e alinhamento com o escopo aprovado.
            \item Facilitei a comunicação em projetos internacionais, atuando como ponte linguística e técnica em reuniões multiculturais para alinhar expectativas e objetivos entre stakeholders globais.
            \item Monitorei logs, métricas e a saúde de pods utilizando Datadog e OpenShift, identificando anomalias e apoiando a análise de incidentes para aumentar a observabilidade e a confiabilidade dos ambientes.
        \end{itemize}

    \cventry{Deskcorp}{Belém, Brasil}{Analista de Testes (QA)}{Maio de 2026 -- Julho de 2026}
        \begin{itemize}
            \item Defini estratégias de testes baseadas em risco, estruturando a pirâmide de testes, os critérios de aceite e os cenários de regressão para priorizar funcionalidades críticas e antecipar falhas no ciclo de desenvolvimento.
            \item Implementei e mantive uma suíte abrangente de testes unitários, de integração e de ponta a ponta utilizando TypeScript, Jest e Testing Library, aplicando mocks, fixtures e dados de teste reutilizáveis para ampliar a cobertura e a confiabilidade das aplicações.
            \item Desenvolvi cenários avançados de testes de performance, incluindo carga, estresse, spike e smoke, utilizando k6 e critérios de aprovação baseados em tempo de resposta, taxa de erros e throughput para identificar gargalos e validar a escalabilidade do sistema.
            \item Mapeei as especificações de negócio e os cenários de validação por meio de uma Matriz de Rastreabilidade de Requisitos e Testes, assegurando a cobertura dos critérios de aceite e das funcionalidades críticas.
            \item Automatizei a esteira de testes utilizando Node.js e integração contínua, configurando quality gates, relatórios de execução e validações de regressão para reduzir o esforço manual e acelerar entregas seguras.
            \item Analisei resultados e indicadores de qualidade, conduzindo a triagem de defeitos, a investigação de causa raiz e o tratamento de testes instáveis para promover melhorias contínuas no processo de desenvolvimento.
        \end{itemize}

    \cventry{Defensoria Pública do Estado do Pará}{Belém, Brasil}{Analista de Sistemas}{Agosto de 2024 -- Maio de 2026}
        \begin{itemize}
            \item Construí aplicações web otimizadas utilizando React, Next.js e Tailwind CSS, implementando renderização híbrida e estilos utilitários que melhoraram a velocidade de carregamento e os índices de Core Web Vitals.
            \item Otimizei o código dos sistemas aplicando os princípios SOLID e a estratégia DRY, eliminando redundâncias e reduzindo em 70\% o tempo de processamento das pipelines.
            \item Colaborei com uma equipe ágil na implementação de uma cobertura de testes robusta utilizando Jest e React Testing Library, melhorando a estabilidade do sistema e reduzindo em 40\% a incidência de bugs em produção.
            \item Coordenei o planejamento de sprints e o acompanhamento de metas em equipes multidisciplinares, utilizando ferramentas ágeis para garantir transparência e alinhamento com os prazos acordados junto aos stakeholders.
            \item Gerenciei o versionamento de código e os processos de code review via GitLab, aumentando a qualidade técnica das entregas e promovendo um ambiente de desenvolvimento mais colaborativo e seguro.
            \item Implementei a conteinerização da aplicação utilizando Docker e Docker Compose, padronizando os ambientes de desenvolvimento e reduzindo o tempo de configuração de novas máquinas.
        \end{itemize}

    \newpage
    \setlist[itemize]{itemsep=0em, topsep=0.1em}

    \cventry{Freelancer}{Remoto}{Desenvolvedor Full Stack}{Junho de 2023 -- Agosto de 2024}
        \begin{itemize}
            \item Implementei sistemas de design escaláveis utilizando React e Shadcn, garantindo consistência visual e facilitando a manutenção de estilos complexos com Styled Components.
            \item Estruturei a arquitetura de serviços e controladores com NestJS, padronizando os endpoints e facilitando a integração entre frontend e backend por meio de uma API bem definida.
            \item Desenvolvi o armazenamento de dados utilizando MongoDB, permitindo um fluxo de informações dinâmico e sem restrições de esquema, com integração fluida entre o backend em Node.js e o banco de dados.
            \item Implementei testes unitários utilizando Jest e React Testing Library, aumentando a confiabilidade dos componentes e facilitando a identificação de falhas durante o desenvolvimento.
            \item Desenvolvi testes de integração e de ponta a ponta utilizando Selenium, validando os principais fluxos da aplicação e reduzindo o risco de regressões entre os sistemas.
            \item Implementei smoke tests para verificar rapidamente as funcionalidades críticas após novas versões, facilitando a identificação antecipada de falhas nas entregas.
            \item Gerenciei a infraestrutura e o deploy da aplicação na plataforma Render, reduzindo em 50\% os custos mensais de hospedagem sem comprometer a performance e a disponibilidade do sistema.
        \end{itemize}

    \cventry{Frexco Alimentos}{Remoto}{Desenvolvedor Full Stack}{Maio de 2022 -- Junho de 2023}
        \begin{itemize}
            \item Arquitetei sistemas web escaláveis utilizando TypeScript, React.js e Material UI, garantindo consistência visual e reutilização de componentes, o que tornou a manutenção mais ágil.
            \item Implementei e mantive serviços independentes utilizando Flask e Django, garantindo uma integração eficiente entre sistemas e reduzindo o tempo de resposta do servidor.
            \item Implementei testes unitários utilizando Jest e React Testing Library, ampliando a cobertura dos componentes e contribuindo para maior estabilidade nas entregas.
            \item Colaborei com times interdisciplinares utilizando Jira e metodologias ágeis, contribuindo para a entrega das sprints dentro dos prazos e para maior transparência no progresso dos projetos.
        \end{itemize}

    \cventry{Cs - Consoft}{Belém, Brasil}{Desenvolvedor Full Stack}{Outubro de 2021 -- Março de 2022}
        \begin{itemize}
            \item Desenvolvi APIs REST utilizando Django REST Framework, garantindo integração eficiente entre sistemas e melhorando a performance das aplicações.
            \item Documentei APIs com DRF-YASG, promovendo clareza e padronização e facilitando a manutenção e a integração dos serviços.
            \item Criei interfaces web responsivas e interativas utilizando Vue.js e Quasar, aplicando boas práticas de experiência do usuário e garantindo consistência visual em diferentes dispositivos.
            \item Colaborei com times interdisciplinares no alinhamento de requisitos e na implementação de soluções, contribuindo para a qualidade das entregas entre as diferentes áreas.
            \item Gerenciei tarefas e sprints ágeis utilizando Trello e Jira, garantindo organização, priorização das demandas e cumprimento dos prazos estabelecidos.
        \end{itemize}

% ------

\section{Habilidades}
    {\footnotesize
    \begin{itemize}[itemsep=0.1em, topsep=0.15em]
        \item \textbf{Desenvolvimento:} TypeScript, JavaScript, Python, React, Vue.js, Next.js, Vite, Vuetify, Quasar, Tailwind CSS, Material UI, Styled Components, Shadcn, Node.js, NestJS, Django, Django REST Framework, Flask, APIs REST, DRF-YASG, MongoDB e MySQL
        \item \textbf{Testes e Qualidade:} Jest, Vitest, Testing Library, React Testing Library, Selenium, k6, testes unitários, de integração, E2E e de performance, smoke tests, automação de testes, quality gates, testes baseados em risco, rastreabilidade e análise de causa raiz
        \item \textbf{Arquitetura e Governança:} Arquitetura de Soluções, modelagem de processos, Bizagi, draw.io, IAM, levantamento de requisitos, gestão de riscos e documentação técnica
        \item \textbf{DevOps e Gestão:} Git, GitHub, GitLab, Docker, Docker Compose, CI/CD, Render, Jira, Trello, OpenProject e metodologias ágeis
        \item \textbf{Idiomas:} Português (Nativo), Inglês (Avançado), Espanhol (Intermediário)
    \end{itemize}
    }

% ------

\section{Educação}
    \cventry{Cesupa}{Belém, Brasil}{Bacharel em Ciências da Computação}{Janeiro de 2019 -- Dezembro de 2022}

\end{document}
